% !TeX program = pdflatex
% UTF-8; standalone source; all diagrams are TikZ.
\documentclass[11pt,a4paper]{article}
\usepackage[utf8]{inputenc}
\usepackage[T5]{fontenc}
\usepackage[vietnamese]{babel}
\usepackage{lmodern}
\usepackage[a4paper,left=19mm,right=19mm,top=21mm,bottom=21mm,headheight=16pt]{geometry}
\usepackage{amsmath,amssymb,booktabs,array,tabularx,longtable,ragged2e}
\usepackage[table]{xcolor}
\usepackage{tikz,listings,enumitem,fancyhdr,lastpage,needspace}
\usetikzlibrary{arrows.meta,positioning,calc}
\usepackage[most]{tcolorbox}
\usepackage[unicode,hidelinks]{hyperref}
\usepackage{bookmark}
\definecolor{Ink}{HTML}{18384F}
\definecolor{Teal}{HTML}{087E8B}
\definecolor{Paper}{HTML}{F1F6F8}
\definecolor{Muted}{HTML}{526477}
\definecolor{Green}{HTML}{276548}
\renewcommand{\tabularxcolumn}[1]{>{\RaggedRight\arraybackslash}p{#1}}
\setlength{\parindent}{0pt}\setlength{\parskip}{5pt}
\setlength{\emergencystretch}{2em}\linespread{1.025}
\setlist[enumerate]{leftmargin=6mm,itemsep=4pt,topsep=3pt}
\setlist[itemize]{leftmargin=5mm,itemsep=3pt,topsep=3pt}
\pagestyle{fancy}\fancyhf{}
\fancyhead[L]{\fontsize{8}{9.5}\selectfont\sffamily\color{Ink}BIỂU DIỄN TRI THỨC VÀ TÌM KIẾM NÂNG CAO}
\fancyhead[R]{\fontsize{8}{9.5}\selectfont\sffamily\VersionShort}
\fancyfoot[L]{\fontsize{8.5}{10}\selectfont\color{Muted}Nguyễn Bích Vân\textbar Biểu diễn tri thức và tìm kiếm nâng cao}
\fancyfoot[R]{\footnotesize\thepage/\pageref*{LastPage}}
\renewcommand{\headrulewidth}{.4pt}
\newtcolorbox{notice}[1][]{enhanced,breakable,colback=Paper,colframe=Teal!45,boxrule=.5pt,arc=1mm,left=3mm,right=3mm,top=2mm,bottom=2mm,#1}
\newcommand{\head}[1]{\par\Needspace{4\baselineskip}\medskip{\large\sffamily\bfseries\color{Ink}#1}\par}
\newcommand{\pagehead}[1]{\clearpage\begin{tcolorbox}[colback=Ink,colframe=Ink,coltext=white,fontupper=\sffamily\large\bfseries,arc=1mm]#1\end{tcolorbox}}
\lstdefinestyle{py}{language=Python,basicstyle=\ttfamily\fontsize{8.5}{10.1}\selectfont,keywordstyle=\bfseries\color{Teal},commentstyle=\color{Muted}\itshape,stringstyle=\color{Green},showstringspaces=false,keepspaces=true,columns=fullflexible,breaklines=true,frame=single,rulecolor=\color{Teal!35},backgroundcolor=\color{Paper},framesep=5pt,aboveskip=6pt,belowskip=6pt,upquote=true,tabsize=4}
\lstset{style=py}
\hypersetup{pdfauthor={Nguyễn Bích Vân},pdflang={vi-VN},pdfsubject={Biểu diễn tri thức và tìm kiếm nâng cao}}
\newcommand{\VersionShort}{NHÓM 1}
\hypersetup{pdftitle={Giữa kỳ 8-puzzle nâng cao -- nhóm 5 -- 10 ngày}}
\begin{document}

{\sffamily\color{Teal}\bfseries BÀI TẬP NHÓM GIỮA KỲ ·}\par\medskip
{\sffamily\color{Ink}\fontsize{26}{30}\selectfont\bfseries Xếp hình 8 ô}\par\smallskip
{\Large\sffamily Heuristic tốt, chi phí đúng, kết quả có thể kiểm chứng}\par\medskip
\textbf{Giảng viên:} TS. Nguyễn Bích Vân \quad IAI--UET--VNU\par
\textbf{Học phần:} Biểu diễn tri thức và tìm kiếm nâng cao\par
\textbf{Nhóm 1: Đoàn Mai Trúc Linh, Nguyễn Thanh Bình, Trần Thùy Dương, Hoàng Thị Ngọc Sơn, Phạm Hồ Minh Thư}  \quad 

\textbf{Hạn nộp bài:} 24-10-2026\par
\textbf{Bảo vệ:} 15--20 phút/nhóm trên lớp. Sau phần trình bày chung của đại diện nhóm, giảng viên có thể hỏi riêng từng thành viên về mức độ hiểu bài. 
\begin{notice}
\textbf{Mã mẫu chạy được ngay.} Được dùng nguyên lõi BFS, UCS, A*; không phải viết lại từ đầu.
Nhóm phải thay dữ liệu, giải thích thuật toán, tính tay, thiết kế kiểm thử và nhận xét có bằng chứng.
%Không đánh giá đóng góp chỉ bằng số dòng code hoặc việc đổi tên biến.
\end{notice}

\head{Sản phẩm cuối}
Một gói ZIP có \textbf{báo cáo PDF trình bày thuật toán và nhận xét kiểm thử}, PDF tính tay,
6--8 slide thuyết trình, các tệp Python chạy được, dữ liệu riêng, kết quả JSON/CSV,
log kiểm thử, README tái tạo và bảng đóng góp. 

\par\begingroup\small\renewcommand{\arraystretch}{1.18}
\begin{tabularx}{\linewidth}{lXr}\toprule
\textbf{Mục} & \textbf{Nội dung} & \textbf{Điểm}\\\midrule
A & Mô hình, A* và lập luận về heuristic & 15\\
B & Tính tay, vết A* và phát hiện sai lầm & 20\\
C & Chạy mã, thay dữ liệu và tạo kết quả & 20\\
D & Kiểm thử và phân tích thực nghiệm & 20\\
E & Báo cáo, khả năng tái tạo và cộng tác & 10\\
F & Bảo vệ riêng từng thành viên & 15\\
\bottomrule\end{tabularx}\endgroup\par

Điểm mỗi sinh viên: $(G+F_i)/10$, với $G\le85$ là điểm sản phẩm chung và $F_i\le15$
là điểm bảo vệ của người $i$. 
\pagehead{1. Luật chơi, chi phí và heuristic}
Một trạng thái là hoán vị của $0,1,\ldots,8$ trên bàn $3\times3$. Số 0 là ô trống.
Một hành động đổi 0 với ô kề \textbf{trên/dưới/trái/phải}; không đi chéo.
Dùng tuple đọc theo hàng. Hàng/cột có chỉ số $0,1,2$.
\begin{center}
\begin{tikzpicture}[x=0.58cm,y=0.58cm,baseline=(current bounding box.center)]
\filldraw[fill=Paper,draw=Ink!55] (0,0) rectangle ++(1,-1);
\node[font=\sffamily\bfseries] at (0.5,-0.5) {1};
\filldraw[fill=Paper,draw=Ink!55] (1,0) rectangle ++(1,-1);
\node[font=\sffamily\bfseries] at (1.5,-0.5) {2};
\filldraw[fill=Paper,draw=Ink!55] (2,0) rectangle ++(1,-1);
\node[font=\sffamily\bfseries] at (2.5,-0.5) {3};
\filldraw[fill=Paper,draw=Ink!55] (0,-1) rectangle ++(1,-1);
\node[font=\sffamily\bfseries] at (0.5,-1.5) {4};
\filldraw[fill=Paper,draw=Ink!55] (1,-1) rectangle ++(1,-1);
\node[font=\sffamily\bfseries] at (1.5,-1.5) {5};
\filldraw[fill=Paper,draw=Ink!55] (2,-1) rectangle ++(1,-1);
\node[font=\sffamily\bfseries] at (2.5,-1.5) {6};
\filldraw[fill=Paper,draw=Ink!55] (0,-2) rectangle ++(1,-1);
\node[font=\sffamily\bfseries] at (0.5,-2.5) {7};
\filldraw[fill=Paper,draw=Ink!55] (1,-2) rectangle ++(1,-1);
\node[font=\sffamily\bfseries] at (1.5,-2.5) {8};
\filldraw[fill=white,draw=Ink!55] (2,-2) rectangle ++(1,-1);
\node[font=\small\sffamily] at (1.5,-3.45) {Đích chuẩn $G$};
\end{tikzpicture}\hspace{20mm}\begin{tikzpicture}[x=0.58cm,y=0.58cm,baseline=(current bounding box.center)]
\filldraw[fill=Paper,draw=Ink!55] (0,0) rectangle ++(1,-1);
\node[font=\sffamily\bfseries] at (0.5,-0.5) {1};
\filldraw[fill=Paper,draw=Ink!55] (1,0) rectangle ++(1,-1);
\node[font=\sffamily\bfseries] at (1.5,-0.5) {2};
\filldraw[fill=Paper,draw=Ink!55] (2,0) rectangle ++(1,-1);
\node[font=\sffamily\bfseries] at (2.5,-0.5) {3};
\filldraw[fill=Paper,draw=Ink!55] (0,-1) rectangle ++(1,-1);
\node[font=\sffamily\bfseries] at (0.5,-1.5) {4};
\filldraw[fill=white,draw=Ink!55] (1,-1) rectangle ++(1,-1);
\filldraw[fill=Paper,draw=Ink!55] (2,-1) rectangle ++(1,-1);
\node[font=\sffamily\bfseries] at (2.5,-1.5) {6};
\filldraw[fill=Paper,draw=Ink!55] (0,-2) rectangle ++(1,-1);
\node[font=\sffamily\bfseries] at (0.5,-2.5) {7};
\filldraw[fill=Paper,draw=Ink!55] (1,-2) rectangle ++(1,-1);
\node[font=\sffamily\bfseries] at (1.5,-2.5) {5};
\filldraw[fill=Paper,draw=Ink!55] (2,-2) rectangle ++(1,-1);
\node[font=\sffamily\bfseries] at (2.5,-2.5) {8};
\node[font=\small\sffamily] at (1.5,-3.45) {Ví dụ làm quen $D$};
\end{tikzpicture}
\end{center}
\textbf{U, D, L, R là hướng của ô trống}, không phải hướng của quân số. Mặc định xét
các con theo U, D, L, R; bỏ hành động không hợp lệ.

\par\begingroup\small\renewcommand{\arraystretch}{1.18}
\begin{tabularx}{\linewidth}{lXX}\toprule
\textbf{Mô hình} & \textbf{Chi phí một bước} & \textbf{Mục tiêu}\\\midrule
\texttt{unit} & U/D/L/R đều bằng 1 & Ít bước nhất = ít chi phí nhất\\
\texttt{directional} & U/D bằng 1; L/R bằng 2 & Ít tổng chi phí nhất, không chỉ ít bước nhất\\
\bottomrule\end{tabularx}\endgroup\par

Với $i=1,\ldots,8$, đặt $(r_i,c_i)$ là vị trí hiện tại và $(r_i^*,c_i^*)$ là vị trí trong
\emph{đích đang được truyền vào}. Không mặc định đích luôn là $G$ khi đã thay dữ liệu.
\[
\begin{aligned}
h_1(s)&=\#\{i\ne0:\text{ quân }i\text{ sai vị trí}\},\\
h_2(s)&=\sum_{i=1}^{8}(|r_i-r_i^*|+|c_i-c_i^*|),\\
h_3(s)&=\sum_{i=1}^{8}(|r_i-r_i^*|+2|c_i-c_i^*|).
\end{aligned}
\]
\textbf{$h_3$ là phần mở rộng của bản mới, dùng trong mô hình directional.}
Trong mô hình unit, tên \texttt{weighted} trong mã tự dùng trọng số ngang 1, trùng $h_2$.
Không cộng khoảng cách của 0 vào bất kỳ heuristic chuẩn nào.
\begin{notice}
A* dùng $f=g+h$, với $g$ là tổng \emph{chi phí cạnh} thật đã đi. UCS dùng $h=0$.
BFS vẫn dùng hàng đợi FIFO; ở directional nó tìm ít bước, không được gắn nhãn ``ít chi phí''.
\end{notice}
\head{Khi nào bài toán có lời giải?}
Cho bàn $n\times n$: bỏ 0, đếm cặp nghịch thế so với thứ tự các quân khác 0 trong đích. Với $n$ lẻ,
bài toán có lời giải khi và chỉ khi số nghịch thế chẵn. Được dùng hàm \texttt{is\_solvable}; không cần
phải chứng minh lại đầy đủ định lý này. 

\pagehead{2. Nhiệm vụ A và quy ước đối chiếu}
\head{A1. Mô hình -- 3 điểm}
Trình bày trạng thái, hành động, hàm chuyển, đích, chi phí và lời giải. Phân biệt
độ dài đường với tổng chi phí; giải thích tại sao không được dùng duy nhất vị trí 0 làm trạng thái.
\head{A2. A* và các thuật toán đối chứng -- 5 điểm}
Viết lại mã giả A*: khởi tạo; chọn nút; bỏ bản ghi lỗi thời; kiểm tra đích;
sinh con; cải thiện $g$; lưu cha và dựng đường. Giải thích vai trò của BFS/UCS khi đối chiếu.
Nêu vì sao kiểm tra đích khi \emph{lấy ra hợp lệ}, không khi vừa sinh. Chỉ rõ các đoạn tương ứng trong \texttt{search.py}.
\head{A3. Heuristic và điều kiện tối ưu -- 7 điểm}
Với mô hình unit, chứng minh $h_1\le h_2\le h^*$ và tính nhất quán của $h_1,h_2$.
Với directional, chứng minh $h_1\le h_2\le h_3\le h^*$ và tính nhất quán của $h_3$.
Giải thích vì sao $h_2$ vẫn dùng được khi chi phí ngang tăng thành 2.
Nêu điều kiện bảo đảm tối ưu của A* trong các lần chạy chuẩn; không suy từ một lần chạy nhanh rằng
một heuristic luôn nhanh hơn theo thời gian thực tế.
\head{Quy ước cài đặt để chấm nhất quán}
\begin{enumerate}[label=\arabic*.]
\item Mọi heuristic chuẩn không âm và bằng 0 tại đích. Hòa $f$: FIFO theo số thứ tự đưa vào heap,
\emph{không} dùng trạng thái dạng tuple làm khóa phụ để phá hòa.
\item BFS đánh dấu phát hiện khi enqueue. UCS/A* lưu \texttt{best\_g}; chỉ thêm bản ghi mới
khi $g$ tốt hơn \textbf{hẳn}. Không đổi cha ở chi phí bằng nhau.
\item Bản ghi cũ còn trong heap đến khi pop. Nếu $g$ đã lỗi thời thì bỏ, không mở rộng.
Mã mẫu hỗ trợ mở lại trạng thái khi chi phí thấp hơn lần mở rộng trước; các heuristic chuẩn ở bài này là nhất quán.
\item \texttt{expanded}: số lần thực sự sinh con, không gồm đích.
\texttt{generated}: số con hợp lệ sinh ra, kể cả trùng. \texttt{pushed}: số bản ghi thêm vào biên, kể cả gốc.
\item \texttt{popped}: mọi lần lấy bản ghi ra, kể cả lỗi thời và đích. \texttt{stale\_pops} đếm lần bỏ bản ghi cũ.
\texttt{max\_frontier}: số bản ghi vật lý lớn nhất trong biên, kể cả bản ghi lỗi thời; không phải dung lượng RAM.
\item Ở lần giải thành công: $\texttt{popped}=\texttt{expanded}+\texttt{stale\_pops}+1$.
Tiền xử lý parity không phải mở rộng nút. Chạm giới hạn tài nguyên trả \texttt{limit}, không được ghi là vô nghiệm.
\end{enumerate}
Có thể khảo sát quy ước khác ở phần bổ sung, nhưng bảng chuẩn phải giữ những quy tắc trên.

\pagehead{3. Nhiệm vụ B -- Tính tay và kiểm chứng (20 điểm)}
Dùng trạng thái $H=$ \texttt{(1,3,6,5,0,2,4,7,8)}, đích $G$ như trang trước.
\begin{center}
\begin{tikzpicture}[x=0.58cm,y=0.58cm,baseline=(current bounding box.center)]
\filldraw[fill=Paper,draw=Ink!55] (0,0) rectangle ++(1,-1);
\node[font=\sffamily\bfseries] at (0.5,-0.5) {1};
\filldraw[fill=Paper,draw=Ink!55] (1,0) rectangle ++(1,-1);
\node[font=\sffamily\bfseries] at (1.5,-0.5) {3};
\filldraw[fill=Paper,draw=Ink!55] (2,0) rectangle ++(1,-1);
\node[font=\sffamily\bfseries] at (2.5,-0.5) {6};
\filldraw[fill=Paper,draw=Ink!55] (0,-1) rectangle ++(1,-1);
\node[font=\sffamily\bfseries] at (0.5,-1.5) {5};
\filldraw[fill=white,draw=Ink!55] (1,-1) rectangle ++(1,-1);
\filldraw[fill=Paper,draw=Ink!55] (2,-1) rectangle ++(1,-1);
\node[font=\sffamily\bfseries] at (2.5,-1.5) {2};
\filldraw[fill=Paper,draw=Ink!55] (0,-2) rectangle ++(1,-1);
\node[font=\sffamily\bfseries] at (0.5,-2.5) {4};
\filldraw[fill=Paper,draw=Ink!55] (1,-2) rectangle ++(1,-1);
\node[font=\sffamily\bfseries] at (1.5,-2.5) {7};
\filldraw[fill=Paper,draw=Ink!55] (2,-2) rectangle ++(1,-1);
\node[font=\sffamily\bfseries] at (2.5,-2.5) {8};
\node[font=\small\sffamily] at (1.5,-3.45) {Trạng thái $H$};
\end{tikzpicture}\hspace{18mm}\begin{tikzpicture}[x=0.58cm,y=0.58cm,baseline=(current bounding box.center)]
\filldraw[fill=Paper,draw=Ink!55] (0,0) rectangle ++(1,-1);
\node[font=\sffamily\bfseries] at (0.5,-0.5) {1};
\filldraw[fill=Paper,draw=Ink!55] (1,0) rectangle ++(1,-1);
\node[font=\sffamily\bfseries] at (1.5,-0.5) {2};
\filldraw[fill=Paper,draw=Ink!55] (2,0) rectangle ++(1,-1);
\node[font=\sffamily\bfseries] at (2.5,-0.5) {3};
\filldraw[fill=Paper,draw=Ink!55] (0,-1) rectangle ++(1,-1);
\node[font=\sffamily\bfseries] at (0.5,-1.5) {4};
\filldraw[fill=Paper,draw=Ink!55] (1,-1) rectangle ++(1,-1);
\node[font=\sffamily\bfseries] at (1.5,-1.5) {5};
\filldraw[fill=white,draw=Ink!55] (2,-1) rectangle ++(1,-1);
\filldraw[fill=Paper,draw=Ink!55] (0,-2) rectangle ++(1,-1);
\node[font=\sffamily\bfseries] at (0.5,-2.5) {7};
\filldraw[fill=Paper,draw=Ink!55] (1,-2) rectangle ++(1,-1);
\node[font=\sffamily\bfseries] at (1.5,-2.5) {8};
\filldraw[fill=Paper,draw=Ink!55] (2,-2) rectangle ++(1,-1);
\node[font=\sffamily\bfseries] at (2.5,-2.5) {6};
\node[font=\small\sffamily] at (1.5,-3.45) {Bàn $V$ cho câu B4};
\end{tikzpicture}
\end{center}
\head{B1 -- 5 điểm: đặc trưng và các con}
Tính $h_1,h_2,h_3$ tại $H$; chỉ rõ đóng góp của từng quân. Sinh các con hợp lệ theo UDLR;
lập bảng trạng thái con, chi phí cạnh unit/directional và ba heuristic.
\head{B2 -- 6 điểm: sáu lần lấy ra đầu của A*}
Dùng unit, $h_2$ và FIFO khi hòa. Thực hiện \textbf{6 lần lấy bản ghi hợp lệ ra} (tính gốc là lần 1).
Ở mỗi lần ghi trạng thái, $g,h,f$ và \textbf{OPEN sau khi xử lý, sắp theo thứ tự sẽ lấy ra}.
Đặt nhãn đường như $S/L/D$ để đỡ vẽ lặp bàn; cung cấp bảng ánh xạ nhãn--trạng thái.
Chỉ đưa bản ghi cải thiện vào OPEN; không đưa lại cha bằng một đường dài hơn.
\head{B3 -- 4 điểm: đường đi và chứng minh tối ưu}
Tìm một đường hợp lệ từ $H$ về $G$; ghi dãy hành động, số bước và chi phí ở cả hai mô hình.
Dùng các cận heuristic vừa tính để chứng minh đường của nhóm tối ưu, không chỉ viết ``vì A* trả ra''.
Chỉ phần B1/B2 phải hoàn thành trước khi mở kết quả code; B3 được dùng code gợi ý đường rồi tự kiểm chứng từng bước.
\head{B4 -- 5 điểm: ba lỗi lập luận}
\begin{enumerate}[label=\alph*)]
\item Một hàm ``số ô sai'' đếm cả 0. Tính hàm đó trên $V$ và chi phí thật về $G$; kết luận
về chấp nhận được và nhất quán. Không chỉ nói chung rằng ``không nên đếm 0''.
\item Đổi đích thành $G'=$ \texttt{(1,2,3,4,5,6,0,7,8)} nhưng vẫn tính Manhattan theo $G$.
Tính hàm sai ngay tại $G'$ và chỉ ra điều kiện bị vi phạm.
\item Với $Z=$ \texttt{(1,2,3,4,5,6,8,7,0)}, tính parity so với $G$.
Phân biệt kết luận vô nghiệm do parity với một lần chạy dừng vì \texttt{limit}.
\end{enumerate}
Nộp PDF tính tay có công thức, bảng và diễn giải, không chỉ ảnh chụp đầu ra Python.

\pagehead{4. Nhiệm vụ C -- Thực hành với mã hoàn chỉnh (20 điểm)}
\head{C1 -- 8 điểm: 36 cấu hình chung}
Giữ nguyên bốn trạng thái trong \texttt{data\_cases.py} và đích $G$:
\begin{center}
\begin{tikzpicture}[x=0.48cm,y=0.48cm,baseline=(current bounding box.center)]
\filldraw[fill=Paper,draw=Ink!55] (0,0) rectangle ++(1,-1);
\node[font=\sffamily\bfseries] at (0.5,-0.5) {1};
\filldraw[fill=Paper,draw=Ink!55] (1,0) rectangle ++(1,-1);
\node[font=\sffamily\bfseries] at (1.5,-0.5) {3};
\filldraw[fill=Paper,draw=Ink!55] (2,0) rectangle ++(1,-1);
\node[font=\sffamily\bfseries] at (2.5,-0.5) {6};
\filldraw[fill=Paper,draw=Ink!55] (0,-1) rectangle ++(1,-1);
\node[font=\sffamily\bfseries] at (0.5,-1.5) {5};
\filldraw[fill=white,draw=Ink!55] (1,-1) rectangle ++(1,-1);
\filldraw[fill=Paper,draw=Ink!55] (2,-1) rectangle ++(1,-1);
\node[font=\sffamily\bfseries] at (2.5,-1.5) {2};
\filldraw[fill=Paper,draw=Ink!55] (0,-2) rectangle ++(1,-1);
\node[font=\sffamily\bfseries] at (0.5,-2.5) {4};
\filldraw[fill=Paper,draw=Ink!55] (1,-2) rectangle ++(1,-1);
\node[font=\sffamily\bfseries] at (1.5,-2.5) {7};
\filldraw[fill=Paper,draw=Ink!55] (2,-2) rectangle ++(1,-1);
\node[font=\sffamily\bfseries] at (2.5,-2.5) {8};
\node[font=\small\sffamily] at (1.5,-3.45) {M1};
\end{tikzpicture}\hspace{5mm}\begin{tikzpicture}[x=0.48cm,y=0.48cm,baseline=(current bounding box.center)]
\filldraw[fill=Paper,draw=Ink!55] (0,0) rectangle ++(1,-1);
\node[font=\sffamily\bfseries] at (0.5,-0.5) {1};
\filldraw[fill=Paper,draw=Ink!55] (1,0) rectangle ++(1,-1);
\node[font=\sffamily\bfseries] at (1.5,-0.5) {3};
\filldraw[fill=Paper,draw=Ink!55] (2,0) rectangle ++(1,-1);
\node[font=\sffamily\bfseries] at (2.5,-0.5) {8};
\filldraw[fill=Paper,draw=Ink!55] (0,-1) rectangle ++(1,-1);
\node[font=\sffamily\bfseries] at (0.5,-1.5) {7};
\filldraw[fill=Paper,draw=Ink!55] (1,-1) rectangle ++(1,-1);
\node[font=\sffamily\bfseries] at (1.5,-1.5) {2};
\filldraw[fill=Paper,draw=Ink!55] (2,-1) rectangle ++(1,-1);
\node[font=\sffamily\bfseries] at (2.5,-1.5) {4};
\filldraw[fill=Paper,draw=Ink!55] (0,-2) rectangle ++(1,-1);
\node[font=\sffamily\bfseries] at (0.5,-2.5) {5};
\filldraw[fill=Paper,draw=Ink!55] (1,-2) rectangle ++(1,-1);
\node[font=\sffamily\bfseries] at (1.5,-2.5) {6};
\filldraw[fill=white,draw=Ink!55] (2,-2) rectangle ++(1,-1);
\node[font=\small\sffamily] at (1.5,-3.45) {M2};
\end{tikzpicture}\hspace{5mm}\begin{tikzpicture}[x=0.48cm,y=0.48cm,baseline=(current bounding box.center)]
\filldraw[fill=white,draw=Ink!55] (0,0) rectangle ++(1,-1);
\filldraw[fill=Paper,draw=Ink!55] (1,0) rectangle ++(1,-1);
\node[font=\sffamily\bfseries] at (1.5,-0.5) {8};
\filldraw[fill=Paper,draw=Ink!55] (2,0) rectangle ++(1,-1);
\node[font=\sffamily\bfseries] at (2.5,-0.5) {2};
\filldraw[fill=Paper,draw=Ink!55] (0,-1) rectangle ++(1,-1);
\node[font=\sffamily\bfseries] at (0.5,-1.5) {4};
\filldraw[fill=Paper,draw=Ink!55] (1,-1) rectangle ++(1,-1);
\node[font=\sffamily\bfseries] at (1.5,-1.5) {6};
\filldraw[fill=Paper,draw=Ink!55] (2,-1) rectangle ++(1,-1);
\node[font=\sffamily\bfseries] at (2.5,-1.5) {3};
\filldraw[fill=Paper,draw=Ink!55] (0,-2) rectangle ++(1,-1);
\node[font=\sffamily\bfseries] at (0.5,-2.5) {7};
\filldraw[fill=Paper,draw=Ink!55] (1,-2) rectangle ++(1,-1);
\node[font=\sffamily\bfseries] at (1.5,-2.5) {1};
\filldraw[fill=Paper,draw=Ink!55] (2,-2) rectangle ++(1,-1);
\node[font=\sffamily\bfseries] at (2.5,-2.5) {5};
\node[font=\small\sffamily] at (1.5,-3.45) {M3};
\end{tikzpicture}\hspace{5mm}\begin{tikzpicture}[x=0.48cm,y=0.48cm,baseline=(current bounding box.center)]
\filldraw[fill=Paper,draw=Ink!55] (0,0) rectangle ++(1,-1);
\node[font=\sffamily\bfseries] at (0.5,-0.5) {7};
\filldraw[fill=Paper,draw=Ink!55] (1,0) rectangle ++(1,-1);
\node[font=\sffamily\bfseries] at (1.5,-0.5) {5};
\filldraw[fill=Paper,draw=Ink!55] (2,0) rectangle ++(1,-1);
\node[font=\sffamily\bfseries] at (2.5,-0.5) {1};
\filldraw[fill=Paper,draw=Ink!55] (0,-1) rectangle ++(1,-1);
\node[font=\sffamily\bfseries] at (0.5,-1.5) {6};
\filldraw[fill=Paper,draw=Ink!55] (1,-1) rectangle ++(1,-1);
\node[font=\sffamily\bfseries] at (1.5,-1.5) {8};
\filldraw[fill=Paper,draw=Ink!55] (2,-1) rectangle ++(1,-1);
\node[font=\sffamily\bfseries] at (2.5,-1.5) {4};
\filldraw[fill=Paper,draw=Ink!55] (0,-2) rectangle ++(1,-1);
\node[font=\sffamily\bfseries] at (0.5,-2.5) {2};
\filldraw[fill=Paper,draw=Ink!55] (1,-2) rectangle ++(1,-1);
\node[font=\sffamily\bfseries] at (1.5,-2.5) {3};
\filldraw[fill=white,draw=Ink!55] (2,-2) rectangle ++(1,-1);
\node[font=\small\sffamily] at (1.5,-3.45) {M4};
\end{tikzpicture}
\end{center}

\par\begingroup\small\renewcommand{\arraystretch}{1.18}
\begin{tabularx}{\linewidth}{lX}\toprule
\textbf{Case} & \textbf{Tuple đọc theo hàng}\\\midrule
M1 & \texttt{(1,3,6,5,0,2,4,7,8)}\\
M2 & \texttt{(1,3,8,7,2,4,5,6,0)}\\
M3 & \texttt{(0,8,2,4,6,3,7,1,5)}\\
M4 & \texttt{(7,5,1,6,8,4,2,3,0)}\\
\bottomrule\end{tabularx}\endgroup\par

\par\begingroup\small\renewcommand{\arraystretch}{1.18}
\begin{tabularx}{\linewidth}{lXl}\toprule
\textbf{Mô hình} & \textbf{Thuật toán/heuristic} & \textbf{Số dòng}\\\midrule
unit & BFS; UCS; A*/$h_1$; A*/$h_2$ & 4 bàn $	imes$ 4 = 16\\
directional & BFS; UCS; A*/$h_1$; A*/$h_2$; A*/$h_3$ & 4 bàn $	imes$ 5 = 20\\
\bottomrule\end{tabularx}\endgroup\par

Chạy \texttt{python run.py --suite --tag bat\_buoc}. Lưu cả JSON (có đường đi) và CSV (số liệu).
Không yêu cầu chép cả 36 dòng vào phần chính báo cáo. Các đường tối ưu có thể khác nhau khi đổi quy tắc hòa.
\head{C2 -- 6 điểm: hai trạng thái riêng, 18 cấu hình}
Sao chép \texttt{du\_lieu\_nhom\_mau.json} thành \texttt{du\_lieu\_nhom.json}; thay cả hai trạng thái.
Hai bàn N1,N2 phải khác nhau, không trùng M1--M4 hoặc ví dụ mẫu, có lời giải tối ưu từ 6 đến 24 bước.
\textbf{Dùng một đích khác $G$ cho cả N1,N2}. Có thể dùng $G'$ ở B4 hoặc một đích khác hợp lệ.
Dùng \texttt{tao\_du\_lieu.py} để sinh từ đích rồi kiểm tra độ dài bằng BFS; số bước xáo trộn không phải đáp số tối ưu.
Chạy hai bàn với cùng chín cấu hình: \texttt{python run.py --data du\_lieu\_nhom.json --suite --tag tu\_tao}.
\head{C3 -- 4 điểm: giải thích một lần chạy}
Trích sáu dòng vết A* trên H để đối chiếu phần B2. Đổi thứ tự con sang \texttt{RDLU},
chạy lại riêng trên H; dự đoán điều có thể thay đổi và kiểm tra cost vẫn tối ưu.
Chỉ đúng dòng mã chọn nút, cập nhật $g$, bỏ bản ghi cũ và dựng đường. Giữ kết quả đổi thứ tự tách khỏi ma trận chuẩn.
\head{C4 -- 2 điểm: kiểm đường và tái chạy}
Kiểm từng bước của mọi đường trả về và tính lại cost bằng luật, không chỉ so số cuối.
Giải nén bài nộp trên máy thành viên khác, chạy đúng README. Ghi rõ tệp nào dùng lại, tệp nào thay/viết.

\pagehead{5. Nhiệm vụ D -- Kiểm thử và nhận xét (20 điểm)}
\head{D1 -- 8 điểm: kiểm thử do nhóm tự thiết kế}
Viết ít nhất \textbf{10 kiểm thử tự động} trong \texttt{test\_nhom.py}; mỗi thành viên phụ trách
ít nhất 2. Ít nhất 3 kỳ vọng phải tính tay độc lập; không lấy đầu ra của chính hàm đang kiểm
làm đáp án của nó. Các test công khai/mẫu được giữ để hồi quy nhưng không thay 10 test của nhóm.
Bao phủ các nhóm tình huống sau (một test có thể có nhiều khẳng định liên quan):

\par\begingroup\small\renewcommand{\arraystretch}{1.18}
\begin{tabularx}{\linewidth}{lX}\toprule
\textbf{Nhóm} & \textbf{Điều phải kiểm}\\\midrule
1--2 & Đầu trùng đích; một bước ngang và một bước dọc với chi phí khác nhau.\\
3--4 & Đích tùy chỉnh; loại 0 khỏi heuristic và kiểm $h(	ext{goal})=0$.\\
5--6 & Vô nghiệm theo parity; dữ liệu sai hoặc hành động vượt biên.\\
7--8 & Tính hợp lệ đường đi/cost; không làm đổi đầu vào.\\
9--10 & A* với $h=0$ đối chiếu UCS; đẳng thức số đếm, trace bật/tắt và giới hạn tài nguyên.\\
\bottomrule\end{tabularx}\endgroup\par

Phải có thêm một khẳng định minh họa \textbf{BFS không bảo đảm chi phí tối ưu ở directional};
có thể gộp vào nhóm 7--10 bằng một bàn đã có. Lưu log thực thi; test bị comment không được tính là đã chạy.
\head{D2 -- 6 điểm: ba phép so sánh có căn cứ}
(1) So $h_1,h_2$ trên unit; (2) so $h_2,h_3$ trên directional;
(3) so số bước và chi phí BFS/UCS trên directional. Dẫn ít nhất hai bàn có số liệu cụ thể.
Với baseline mở rộng $N_0>0$, tỷ lệ giảm mở rộng là $1-N/N_0$; giữ nguyên start, goal,
chi phí và thứ tự khi so heuristic. Không so hai mô hình khác nhau như một thí nghiệm kiểm soát.
\head{D3 -- 4 điểm: bắt lỗi bằng phản ví dụ}
Từ B4, tự viết hai hàm heuristic sai trong \texttt{test\_nhom.py} để minh họa:
đếm cả ô trống và dùng vị trí đích cố định. Viết test \emph{xác nhận lỗi được phát hiện},
không để một assert sai làm toàn bộ bài nộp thất bại. Không sửa heuristic chuẩn thành phiên bản sai.
\head{D4 -- 2 điểm: giới hạn kết luận}
Nêu ít nhất hai giới hạn của thí nghiệm. Phân biệt số bản ghi biên với RAM, số nút với thời gian,
``nhất quán trên mẫu đã thử'' với một chứng minh tổng quát. Không bắt buộc đồ thị hoặc đo thời gian;
nếu đo, dùng cùng máy, tắt trace/step, lặp 3 lần và lấy trung vị.
\begin{notice}
Có mã mẫu không làm mất phần đánh giá hiểu bài. Một đầu ra đúng mà không giải thích
được nguồn dữ liệu, quy tắc cost, heuristic hoặc điều kiện kiểm thử thì chưa đủ điểm các mục tương ứng.
\end{notice}

\pagehead{6. Báo cáo, gói nộp và bảo vệ}
\head{E -- 10 điểm: báo cáo và minh chứng}
Điền mẫu Word hoặc pdfLaTeX trong \texttt{MauBaoCao/}; xuất \texttt{BaoCao.pdf}.
Báo cáo chính \textbf{6--8 trang}, không tính bìa riêng/phụ lục tính tay. Phải có:
trình bày mô hình và A*; giải thích heuristic; đối chiếu tính tay; thiết kế dữ liệu;
bảng so sánh; nhận xét kiểm thử; giới hạn, nguồn dùng lại và đóng góp từng người.
Không nộp mẫu còn nguyên toàn bộ chữ [Điền]; không dán hàng trăm dòng code/trace vào báo cáo.

\textbf{E1 (4 điểm):} giải thích rõ, dẫn bảng/kiểm thử cụ thể. \textbf{E2 (3 điểm):} README tái chạy được,
đủ minh chứng và mã nguồn. \textbf{E3 (3 điểm):} đóng góp kỹ thuật và kiểm tra chéo của cả năm thành viên.
\head{Một gói ZIP}
\begin{lstlisting}[language={}]
Nhom1_8Puzzle_GiuaKy.zip
  Nhom1_8Puzzle/
    BaoCao.pdf           TinhTay.pdf
    ThuyetTrinh.pptx     # hoac ThuyetTrinh.pdf
    README_NHOM.md       DongGop.csv
    puzzle.py           search.py          run.py
    data_cases.py        test_nhom.py
    du_lieu_nhom.json    # them helper .py da su dung
    output/
      bat_buoc.csv       bat_buoc.json      # 36 dong
      tu_tao.csv         tu_tao.json        # 18 dong
      ket_qua_test.txt   # log that; them trace/doi thu tu
\end{lstlisting}
Được nộp thêm tệp hỗ trợ, không nộp môi trường Anaconda, cache, dữ liệu cá nhân không cần thiết. Dùng \texttt{dong\_goi.py --group Nhom1} để tạo ZIP sau khi đủ tệp.
Công cụ chỉ kiểm cấu trúc và số dòng, không chứng nhận báo cáo đúng.
\head{F -- 15 điểm/người: bảo vệ 15--20 phút}
 5 phút đại diện nhóm trình bày; 3 phút chạy minh họa; 8-10 phút hỏi năm thành viên,
khoảng 1.5-2 phút/người. 
Tất cả thành viên phải trả lời; giảng viên được hỏi bất kỳ mô-đun nào.
Năm tiêu chí, mỗi tiêu chí 0--3: mô hình/cost; dự đoán; giải thích dòng mã;
xử lý thay đổi dữ liệu; nhận xét kiểm thử/giới hạn. Ví dụ đổi đích, đổi thứ tự con,
đổi phí ngang, chỉ lỗi h đếm 0. Không chỉ bấm Run và đọc kết quả.
\textbf{Demo phải dùng dữ liệu của nhóm}, có dự đoán trước khi chạy. Không bắt buộc giao diện đồ họa.


\head{Cách dùng mã mẫu và công cụ hỗ trợ}
Được giữ nguyên các hàm thuật toán, runner và luật trong gói. Ghi nguồn và những thay đổi thực sự.
Không gọi một thư viện ngoài để giải thay toàn bộ bài. Báo cáo cần phân biệt phần dùng lại,
phần tự xây dựng dữ liệu, tính tay, kiểm thử và nhận xét. 


\end{document}
